\documentclass[10pt,a4paper]{article}
\usepackage[utf8]{inputenc}
\usepackage[T1]{fontenc}
\usepackage[ngerman]{babel}
\usepackage[margin=1.7cm,top=1.5cm,bottom=1.7cm]{geometry}
\usepackage{titlesec}
\titleformat*{\section}{\normalfont\large\bfseries}
\titleformat*{\subsection}{\normalfont\normalsize\bfseries}
\titlespacing*{\section}{0pt}{8pt}{3pt}
\titlespacing*{\subsection}{0pt}{6pt}{2pt}
\usepackage{amsmath,amssymb}
\usepackage{listings}
\usepackage{xcolor}
\usepackage{fancyhdr}
\usepackage{parskip}
\usepackage{needspace}
\usepackage{tikz}
\usepackage{array}

\newcommand{\mcbox}{\raisebox{-1pt}{\fbox{\rule{0pt}{7pt}\rule{7pt}{0pt}}}\hspace{6pt}}
\newcommand{\antwortlinien}[1]{%
  \par\vspace{2pt}%
  \newcount\zeilen \zeilen=#1
  \loop
    \noindent\rule{\linewidth}{0.4pt}\par\vspace{11pt}%
    \advance\zeilen by -1
  \ifnum\zeilen>0 \repeat
}
\newcommand{\pkt}[1]{\hfill\textbf{(#1 P)}}
\lstset{language=Python,basicstyle=\ttfamily\footnotesize,keywordstyle=\bfseries,
  commentstyle=\itshape\color{gray},showstringspaces=false,frame=single,
  framerule=0.3pt,xleftmargin=8pt,framexleftmargin=6pt,aboveskip=6pt,belowskip=6pt}

\pagestyle{fancy}
\fancyhf{}
\fancyhead[L]{\small MDT5/2 -- Probeklausur 2 (generiert)}
\fancyhead[R]{\small Maschinelles Lernen zur Anomalieerkennung}
\fancyfoot[C]{\small Seite \thepage}
\renewcommand{\headrulewidth}{0.4pt}

\begin{document}

\begin{center}
  {\Large\bfseries Probeklausur 2}\\[4pt]
  {\large Maschinelles Lernen zur Anomalieerkennung -- Kapitel 1--8}\\[6pt]
  \begin{tabular}{@{}ll@{}}
  \textbf{Bearbeitungszeit:} & 120 Minuten (bewusst mehr als die echten 90)\\
  \textbf{Hilfsmittel:} & Skript (ausgedruckt) + handgeschriebene Blätter. \textbf{Keine Praktika.}\\
  \textbf{Gesamtpunktzahl:} & 125 Punkte\\
  \end{tabular}
\end{center}

\vspace{2pt}
\noindent\rule{\linewidth}{0.8pt}
\emph{Hinweis zur Bewertung von Aufgabe 1: Für ein richtiges Kreuz \textbf{mit} passender Begründung
gibt es 3 Punkte, für eine richtige Begründung \textbf{allein} 1 Punkt -- auch wenn das Kreuz falsch
gesetzt wurde. Ein Kreuz ohne Begründung gibt 0 Punkte.}

\noindent\rule{\linewidth}{0.8pt}

\section*{Aufgabe 1 -- Multiple Choice mit Begründung \pkt{21}}
\emph{Markieren Sie jeweils die einzige richtige von drei Aussagen und begründen Sie Ihre
Entscheidung in ein bis zwei Sätzen.}

\subsection*{1.1 \ Für die Normalisierung von Merkmalen gilt:}
\mcbox Eine z-Transformation (\texttt{StandardScaler}) bildet die Merkmale \textbf{immer} auf den
Bereich $[-1, 1]$ ab.

\mcbox Für den Isolation Forest ist eine Normalisierung \textbf{essentiell}, weil er die Abstände
zwischen den Punkten auswertet.

\mcbox Für abstandsbasierte Verfahren ist eine Normalisierung nötig, weil sonst das Merkmal mit
dem größten Wertebereich den Abstand dominiert.

Begründung:
\antwortlinien{2}

\subsection*{1.2 \ Für den Parameter \texttt{contamination} der Elliptic Envelope gilt:}
\mcbox \texttt{contamination} gibt die Anzahl der erwarteten Ausreißer im Datensatz an und
verändert den Kovarianz-Fit.

\mcbox \texttt{contamination} ist ein Anteil und wirkt als Schwellwert; die AUC des Modells
ändert sich dadurch \textbf{nicht}.

\mcbox \texttt{contamination} ist ein Formparameter der geschätzten Ellipse und lässt sich
unüberwacht per Grid Search bestimmen.

Begründung:
\antwortlinien{2}

\subsection*{1.3 \ Für einen Autoencoder zur Anomalieerkennung gilt:}
\mcbox Die Sigmoid-Aktivierung in der Ausgabeschicht zeigt an, dass der Autoencoder ein binäres
Klassifikationsproblem (normal / anomal) löst.

\mcbox Hat die Code-Schicht $z$ mindestens so viele Dimensionen wie die Eingabe, kann der
Autoencoder die Identität lernen und wird zur Anomalieerkennung unbrauchbar.

\mcbox Der Autoencoder wird auf Normaldaten \textbf{und} Anomalien trainiert, damit er beide
Klassen unterscheiden kann.

Begründung:
\antwortlinien{2}

\subsection*{1.4 \ Für Gradienten in tiefen Netzen gilt:}
\mcbox ReLU verhindert sowohl verschwindende als auch explodierende Gradienten, weil ihre
Ableitung \textbf{nie} größer als 1 wird.

\mcbox Die Ableitung der Sigmoidfunktion ist höchstens $0{,}25$; über viele Schichten hinweg
verschwindet der Gradient deshalb in Richtung der \textbf{Eingangsschicht}.

\mcbox Explodierende Gradienten treten \textbf{immer} zuerst in der Ausgangsschicht auf, weil
dort der Fehler entsteht.

Begründung:
\antwortlinien{2}

\subsection*{1.5 \ Für die One-Class SVM gilt:}
\mcbox $\nu$ ist ein Strafgewicht für Fehlklassifikationen (wie $C$ bei der Zweiklassen-SVM):
je größer $\nu$, desto härter die Grenze.

\mcbox $\nu$ ist eine obere Schranke für den Anteil der Ausreißer im Training und zugleich eine
untere Schranke für den Anteil der Support-Vektoren.

\mcbox Der Kernel-Trick ist bei der One-Class SVM nur eine Beschleunigung -- man könnte die
Transformation $\varphi(x)$ \textbf{immer} auch explizit ausrechnen.

Begründung:
\antwortlinien{2}

\subsection*{1.6 \ Ein Detektor flaggt in einem Datensatz mit 2\,\% Anomalien sehr viele Punkte.
Dazu gilt:}
\mcbox Die \emph{balanced accuracy} ist hier die ehrlichste Metrik, weil sie die Unbalanciertheit
der Klassen einbezieht.

\mcbox Precision und $F_1$ decken die Fehlalarme auf, weil in ihnen \textbf{keine} True Negatives
vorkommen.

\mcbox Ein hoher Recall belegt, dass nur wenige Fehlalarme auftreten.

Begründung:
\antwortlinien{2}

\subsection*{1.7 \ Für Deep SVDD und GOAD gilt:}
\mcbox Batch Normalization ist im Deep-SVDD-Netz \textbf{grundsätzlich verboten}, weil der
gelernte Parameter $\beta$ wie ein Bias wirkt.

\mcbox $\nu$ legt bei Deep SVDD den Radius als $(1-\nu)$-Quantil der Abstände zu $c$ fest; es
verschiebt nur den Betriebspunkt, die AUC bleibt gleich.

\mcbox Der dritte Term der GOAD-Verlustfunktion ist ein Weight Decay auf die Gewichte des Netzes.

Begründung:
\antwortlinien{2}

\newpage
\section*{Aufgabe 2 -- Ausgabe-Shapes eines Faltungs-Autoencoders \pkt{12}}

Gegeben ist das folgende Netz. Die Eingabe ist ein RGB-Bild der Größe $64\times64\times3$.
Tragen Sie für \textbf{jede} Zeile die Ausgabe-Shape ein (ohne Batch-Dimension).

\begin{lstlisting}
ae = tf.keras.Sequential([
    tf.keras.layers.Input((64, 64, 3)),                                   #  ( 64, 64, 3 )
 1  tf.keras.layers.Conv2D(16, (3,3), padding='same', activation='relu'),
 2  tf.keras.layers.MaxPooling2D((2,2)),
 3  tf.keras.layers.Conv2D(32, (5,5), padding='valid', activation='relu'),
 4  tf.keras.layers.Conv2D(32, (3,3), strides=2, padding='same', activation='relu'),
 5  tf.keras.layers.MaxPooling2D((3,3)),
 6  tf.keras.layers.Flatten(),
 7  tf.keras.layers.Dense(20, activation='relu'),
 8  tf.keras.layers.Dense(512, activation='relu'),
 9  tf.keras.layers.Reshape((4, 4, 32)),
10  tf.keras.layers.Conv2DTranspose(32, (3,3), strides=2, padding='same', activation='relu'),
11  tf.keras.layers.Conv2DTranspose(16, (3,3), strides=2, padding='same', activation='relu'),
12  tf.keras.layers.Conv2DTranspose(8,  (3,3), strides=2, padding='same', activation='relu'),
13  tf.keras.layers.Conv2DTranspose(3,  (3,3), strides=2, padding='same', activation='sigmoid'),
])
\end{lstlisting}

\begin{center}
\begin{tabular}{|c|p{4.2cm}||c|p{4.2cm}|}
\hline
Zeile & Ausgabe-Shape & Zeile & Ausgabe-Shape \\ \hline
1 & \rule{0pt}{13pt} & 8 & \rule{0pt}{13pt}\\ \hline
2 & \rule{0pt}{13pt} & 9 & \rule{0pt}{13pt}\\ \hline
3 & \rule{0pt}{13pt} & 10 & \rule{0pt}{13pt}\\ \hline
4 & \rule{0pt}{13pt} & 11 & \rule{0pt}{13pt}\\ \hline
5 & \rule{0pt}{13pt} & 12 & \rule{0pt}{13pt}\\ \hline
6 & \rule{0pt}{13pt} & 13 & \rule{0pt}{13pt}\\ \hline
7 & \rule{0pt}{13pt} & & \\ \hline
\end{tabular}
\end{center}

\textbf{b)} Wie viele trainierbare Parameter hat Zeile 1, wie viele Zeile 7? Rechenweg angeben.
\antwortlinien{3}

\textbf{c)} Welche Schicht des Netzes trägt das größte Overfitting-Risiko und warum?
\antwortlinien{2}

\textbf{d)} Begründen Sie \textbf{beide} Entwurfsentscheidungen in Zeile 13: warum genau 3 Filter
und warum \texttt{sigmoid}?
\antwortlinien{3}

\newpage
\section*{Aufgabe 3 -- Fehler finden \pkt{18}}

\emph{In den folgenden Ausschnitten können mehrere Fehler stecken -- oder auch keiner.
Benennen Sie jeweils die betroffene Codezeile, den Fehler und seine Auswirkung.}

\subsection*{3.1 \ Training und Evaluierung einer One-Class SVM (Novelty Detection) \pkt{10}}

\begin{lstlisting}
import numpy as np
from sklearn.preprocessing import StandardScaler
from sklearn.model_selection import GridSearchCV, train_test_split
from sklearn.svm import OneClassSVM
from sklearn.metrics import accuracy_score, roc_auc_score

X, y = load_sensor_data()        # y: 0 = normal, 1 = Anomalie; ca. 3 % Anomalien

scaler = StandardScaler()                                                    # (1)
X = scaler.fit_transform(X)                                                  # (2)

X_train, X_test, y_train, y_test = train_test_split(X, y, test_size=0.3)     # (3)

param_grid = {'nu':    np.logspace(-3, 0, 10),                               # (4)
              'gamma': np.logspace(-2, 2, 10)}                               # (5)
search = GridSearchCV(OneClassSVM(kernel='rbf'), param_grid,
                      scoring='accuracy', cv=5)                              # (6)
search.fit(X_train, y_train)                                                 # (7)

best = search.best_estimator_
pred = best.predict(X_train)                                                 # (8)
pred = np.where(pred == 1, 0, 1)
print('Accuracy:', accuracy_score(y_train, pred))                            # (9)
print('AUC:', roc_auc_score(y_test, best.score_samples(X_test)))             # (10)
\end{lstlisting}

\antwortlinien{10}

\newpage
\subsection*{3.2 \ Aufbau eines CNN-Klassifikators für $32\times32\times3$-Bilder, 5 Klassen \pkt{8}}

\begin{lstlisting}
net = tf.keras.Sequential([
    tf.keras.layers.Input((32, 32, 3)),
    tf.keras.layers.Conv2D(16, (3,3), padding='same', activation='relu'),
    tf.keras.layers.MaxPooling2D((2,2)),
    tf.keras.layers.Conv2D(32, (3,3), padding='same', activation='relu'),
    tf.keras.layers.MaxPooling2D((2,2)),
    tf.keras.layers.Conv2D(64, (3,3), padding='same', activation='relu'),
    tf.keras.layers.MaxPooling2D((2,2)),
    tf.keras.layers.MaxPooling2D((4,4)),
    tf.keras.layers.Dense(64, activation='relu'),
    tf.keras.layers.Dense(5,  activation='relu'),
])
net.compile(optimizer='adam', loss='binary_crossentropy', metrics=['accuracy'])
net.fit(x_train, y_train, epochs=20, validation_data=(x_test, y_test))
\end{lstlisting}

\antwortlinien{8}

\newpage
\section*{Aufgabe 4 -- Programmieren \pkt{16}}

\subsection*{4.1 \ Deep-SVDD-Netz in Keras \pkt{10}}

Schreiben Sie das Netz $\varphi(x)$ für eine Deep SVDD auf $28\times28\times1$-Bildern:
zwei Faltungsschichten (die erste mit 8, die zweite mit 16 Feature Maps, jeweils $5\times5$,
\texttt{padding='same'}), auf jede Faltung folgt ein Max-Pooling $2\times2$; danach wird flach
gemacht und über eine vollverknüpfte Schicht auf einen Latent Space der Dimension 32 abgebildet.
Alle Schichten mit L2-Regularisierung ($\lambda = 10^{-6}$), als Aktivierung
\texttt{LeakyReLU}. \textbf{Beachten Sie die Konstruktionsregeln für Deep SVDD} und schreiben Sie
die Shape hinter jede Schicht.

\antwortlinien{12}

\textbf{b)} Schreiben Sie die Zeile, mit der das Zentrum $c$ bestimmt wird. An welcher Stelle des
Programms muss sie stehen, und was passiert, wenn man sie in die Trainingsschleife setzt oder
$c = \mathbf{0}$ wählt?
\antwortlinien{4}

\newpage
\subsection*{4.2 \ Anomalie-Score und Schwellwert in NumPy \pkt{6}}

Gegeben sind ein fertig trainierter Autoencoder \texttt{ae} und Testdaten \texttt{X\_test} der
Form \texttt{(N, 28, 28, 1)}. Aus dem Prozess ist bekannt, dass etwa \textbf{2\,\%} der Messungen
Ausreißer sind. Schreiben Sie den Code, der

\begin{enumerate}
\itemsep0pt
\item den Rekonstruktionsfehler je Sample berechnet,
\item daraus den Schwellwert bestimmt,
\item ein Ergebnis-Array in der sklearn-Konvention ($+1$ = normal, $-1$ = Anomalie) erzeugt,
\item die Anzahl der ausgelösten Alarme ausgibt.
\end{enumerate}

\antwortlinien{8}

\textbf{b)} Über welche Achsen wird beim Rekonstruktionsfehler gemittelt und warum?
Welches Perzentil haben Sie in Schritt 2 verwendet und warum?
\antwortlinien{3}

\newpage
\section*{Aufgabe 5 -- Auswertung eines Score-Histogramms \pkt{14}}

Ein Autoencoder-Ensemble wurde \textbf{unüberwacht} auf $N = 5000$ Messungen angewendet. Das
Histogramm zeigt die Verteilung der Anomalie-Scores $s$ über alle 5000 Punkte
(Klassenbreite 0,5):

\begin{center}
\begin{tikzpicture}[x=1.45cm,y=0.0026cm]
  % Balken
  \fill[gray!35,draw=black,line width=0.3pt] (0,0)   rectangle (0.5,900);
  \fill[gray!35,draw=black,line width=0.3pt] (0.5,0) rectangle (1.0,1500);
  \fill[gray!35,draw=black,line width=0.3pt] (1.0,0) rectangle (1.5,1200);
  \fill[gray!35,draw=black,line width=0.3pt] (1.5,0) rectangle (2.0,700);
  \fill[gray!35,draw=black,line width=0.3pt] (2.0,0) rectangle (2.5,350);
  \fill[gray!35,draw=black,line width=0.3pt] (2.5,0) rectangle (3.0,180);
  \fill[gray!35,draw=black,line width=0.3pt] (3.0,0) rectangle (3.5,90);
  \fill[gray!35,draw=black,line width=0.3pt] (3.5,0) rectangle (4.0,45);
  \fill[gray!35,draw=black,line width=0.3pt] (4.0,0) rectangle (4.5,20);
  \fill[gray!35,draw=black,line width=0.3pt] (4.5,0) rectangle (5.0,10);
  \fill[gray!35,draw=black,line width=0.3pt] (5.0,0) rectangle (5.5,4);
  \fill[gray!35,draw=black,line width=0.3pt] (5.5,0) rectangle (6.0,1);
  % Achsen
  \draw[->,line width=0.5pt] (0,0) -- (6.6,0) node[right,scale=0.85]{$s$};
  \draw[->,line width=0.5pt] (0,0) -- (0,1700) node[above,scale=0.85]{Anzahl};
  \foreach \x in {0,1,2,3,4,5,6}
     \draw (\x,0) -- (\x,-35) node[below,scale=0.75]{\x};
  \foreach \y in {500,1000,1500}
     \draw (0,\y) -- (-0.07,\y) node[left,scale=0.75]{\y};
\end{tikzpicture}
\end{center}

\begin{center}\small
\begin{tabular}{|l|r|r|r|r|r|r|r|r|r|r|r|r|}
\hline
Klasse & 0--0,5 & 0,5--1 & 1--1,5 & 1,5--2 & 2--2,5 & 2,5--3 & 3--3,5 & 3,5--4 & 4--4,5 & 4,5--5 & 5--5,5 & 5,5--6\\ \hline
Anzahl & 900 & 1500 & 1200 & 700 & 350 & 180 & 90 & 45 & 20 & 10 & 4 & 1\\ \hline
\end{tabular}
\end{center}

Nachträglich wird ein \textbf{gelabelter} Validierungssatz verfügbar. Damit ergeben sich für
einige Schwellwerte folgende Betriebspunkte:

\begin{center}\small
\begin{tabular}{|c|c|c|c|c|c|}
\hline
Schwellwert $s_0$ & 1,0 & 2,0 & 2,5 & 3,5 & 4,0\\ \hline
TPR & 0,99 & 0,92 & 0,86 & 0,21 & 0,06\\ \hline
FPR & 0,62 & 0,14 & 0,05 & 0,013 & 0,007\\ \hline
\end{tabular}
\hspace{6pt}
\begin{tabular}{|l|c|}
\hline
AUC des Modells & 0,95\\ \hline
\end{tabular}
\end{center}

\textbf{a)} Sie kennen zunächst \textbf{nur} das Histogramm. Wo würden Sie den Schwellwert legen,
und wie viele Alarme löst Ihre Wahl aus? Rechnen Sie die Anzahl aus der Tabelle aus und prüfen
Sie sie auf Plausibilität.
\antwortlinien{4}

\textbf{b)} Eine Kollegin setzt den Schwellwert stattdessen auf das 98{,}4\,\%-Quantil, weil aus
dem Prozess ein Ausreißeranteil von 1,6\,\% bekannt ist. Welcher Schwellwert ergibt sich aus der
Tabelle, wie viele Punkte werden geflaggt -- und welche Eigenschaft garantiert dieses Vorgehen,
welche nicht?
\antwortlinien{4}

\textbf{c)} Bei diesem Schwellwert werden nur 21\,\% der echten Ausreißer gefunden, die AUC
beträgt aber 0,95. Ist das Modell schlecht oder der Schwellwert? Begründen Sie, und nennen Sie
die Eigenschaft der Score-Verteilung, die diesen Effekt erzeugt.
\antwortlinien{4}

\textbf{d)} Bestimmen Sie aus der Betriebspunkt-Tabelle den besten Schwellwert im Sinne des
ROC-Knies. Geben Sie das Rechenkriterium an und vergleichen Sie den gefundenen Punkt mit
$s_0 = 1{,}0$: was kostet der letzte Rest TPR, in Fehlalarmen ausgedrückt?
\antwortlinien{5}

\textbf{e)} Warum ist das Histogramm für die Wahl des Schwellwertes die schlechteste der drei
möglichen Darstellungen? Welche zwei Darstellungen sind besser und was liest man dort jeweils ab?
\antwortlinien{4}

\newpage
\section*{Aufgabe 6 -- SVM und One-Class SVM \pkt{22}}

\subsection*{6.1 \ Zweiklassen-SVM von Hand \pkt{9}}

Gegeben sind acht Punkte in zwei Klassen:

\begin{center}\small
\begin{tabular}{|l|l|}
\hline
Klasse $\bullet$ ($y = -1$) & $(1|2)$, $(2|1)$, $(2|3)$, $(3|1)$\\ \hline
Klasse $\times$ ($y = +1$) & $(3|4)$, $(5|4)$, $(6|3)$, $(6|5)$, $(7|4)$\\ \hline
\end{tabular}
\end{center}

\begin{center}
\begin{tikzpicture}[scale=0.78]
  \draw[step=1,gray!30,line width=0.3pt] (0,0) grid (8,6);
  \draw[->] (0,0) -- (8.4,0) node[right,scale=0.85]{$x_1$};
  \draw[->] (0,0) -- (0,6.4) node[above,scale=0.85]{$x_2$};
  \foreach \x in {1,...,8} \draw (\x,0) -- (\x,-0.12) node[below,scale=0.7]{\x};
  \foreach \y in {1,...,6} \draw (0,\y) -- (-0.12,\y) node[left,scale=0.7]{\y};
  \foreach \p in {(1,2),(2,1),(2,3),(3,1)}
     \fill \p circle (0.11);
  \foreach \p in {(3,4),(5,4),(6,3),(6,5),(7,4)}
     \draw[line width=1pt] \p ++(-0.13,-0.13) -- ++(0.26,0.26) \p ++(-0.13,0.13) -- ++(0.26,-0.26);
\end{tikzpicture}
\end{center}

\textbf{a)} Zeichnen Sie die Trenngerade für ein \textbf{sehr großes} $C$ ein, markieren Sie alle
Support-Vektoren und geben Sie die Breite der Margin an (Rechenweg).
\antwortlinien{4}

\textbf{b)} $C$ wird nun schrittweise \textbf{kleiner}. Beschreiben Sie, was mit der Lage der
Trenngeraden, der Breite der Margin, der Anzahl der Support-Vektoren und den
Fehlklassifikationen passiert -- und was im Grenzfall $C \to 0$ übrig bleibt.
\antwortlinien{5}

\textbf{c)} Ein neuer Punkt $(4|3)$ trifft ein. Zu welcher Klasse gehört er bei sehr großem $C$?
Begründen Sie \textbf{ohne} zu schätzen, und sagen Sie, ob der Punkt innerhalb oder außerhalb der
Margin liegt.
\antwortlinien{3}

\subsection*{6.2 \ Kernel-Trick \pkt{6}}

Gegeben ist die Transformation
$\varphi(\boldsymbol{x}) = \left(x_1^2,\ \sqrt{2}\,x_1x_2,\ x_2^2\right)$
sowie die beiden Punkte $\boldsymbol{a} = (1|2)$ und $\boldsymbol{b} = (3|1)$.

\textbf{a)} Berechnen Sie $\varphi(\boldsymbol{a}) \cdot \varphi(\boldsymbol{b})$ und
$k(\boldsymbol{a}, \boldsymbol{b}) = (\boldsymbol{a} \cdot \boldsymbol{b})^2$ getrennt voneinander
und vergleichen Sie.
\antwortlinien{4}

\textbf{b)} \textbf{Welche} Rechnungen der SVM kommen mit Skalarprodukten aus -- nennen Sie beide.
Worin unterscheidet sich der RBF-Kernel rechnerisch vom polynomialen, und warum ist der
Kernel-Trick dort nicht nur billiger, sondern die einzige Möglichkeit?
\antwortlinien{5}

\subsection*{6.3 \ One-Class SVM \pkt{7}}

Die Normaldaten liegen als kompakte Wolke rechts oben, der Ursprung ist eingezeichnet:

\begin{center}
\begin{tikzpicture}[scale=0.75]
  \draw[step=1,gray!30,line width=0.3pt] (0,0) grid (8,7);
  \draw[->] (0,0) -- (8.4,0) node[right,scale=0.85]{$x_1$};
  \draw[->] (0,0) -- (0,7.4) node[above,scale=0.85]{$x_2$};
  \foreach \x in {1,...,8} \draw (\x,0) -- (\x,-0.12) node[below,scale=0.7]{\x};
  \foreach \y in {1,...,7} \draw (0,\y) -- (-0.12,\y) node[left,scale=0.7]{\y};
  \foreach \p in {(4,4),(4.6,5),(5,4.3),(5.5,5.2),(4.2,5.6),(5.8,4.1),(5.2,5.8),(4.8,4.8),
                  (6.2,5.0),(3.8,4.9),(5.6,3.9),(4.4,4.2),(6.0,5.7),(5.0,5.3),(4.9,6.1)}
     \fill \p circle (0.1);
  \fill[black] (2.6,6.4) circle (0.1);
  \node[scale=0.7] at (2.6,6.85) {Randpunkt};
\end{tikzpicture}
\end{center}

\textbf{a)} Zeichnen Sie die Trennebene für einen \textbf{linearen} Kernel mit $\nu = 0{,}05$ ein
und markieren Sie deutlich, auf welcher Seite Punkte als \textbf{anomal} gelten. Erklären Sie in
einem Satz, wogegen die OCSVM die Normaldaten hier abgrenzt.
\antwortlinien{3}

\textbf{b)} Was ändert sich an Ihrer Zeichnung, wenn $\nu$ deutlich \textbf{größer} gewählt wird?
Nennen Sie zusätzlich die Folge für Over- bzw. Underfitting bei zu kleinem und zu großem $\nu$.
\antwortlinien{3}

\textbf{c)} Nun wird auf einen RBF-Kernel gewechselt. Beschreiben Sie die Entscheidungsgrenze
(nicht nur die Glockenfunktion!) und was mit ihr passiert, wenn $\gamma$ sehr groß wird. Welchen
Wert nimmt die balanced accuracy im Extremfall an und warum?
\antwortlinien{4}

\newpage
\section*{Aufgabe 7 -- Verfahrenswahl bei mehreren Clustern \pkt{12}}

Die folgende Abbildung zeigt die \textbf{Normaldaten} eines Prozesses (zwei Merkmale, ohne
Anomalien -- Novelty-Setup). Zusätzlich sind zwei neue Messungen $\times$ eingezeichnet, die
bewertet werden sollen.

\begin{center}
\begin{tikzpicture}[scale=0.8]
  \draw[step=1,gray!25,line width=0.3pt] (0,0) grid (10,10);
  \draw[->] (0,0) -- (10.4,0) node[right,scale=0.85]{$x_1$};
  \draw[->] (0,0) -- (0,10.4) node[above,scale=0.85]{$x_2$};
  \foreach \x in {2,4,6,8,10} \draw (\x,0) -- (\x,-0.12) node[below,scale=0.7]{\x};
  \foreach \y in {2,4,6,8,10} \draw (0,\y) -- (-0.12,\y) node[left,scale=0.7]{\y};
  % Cluster 1
  \foreach \p in {(1.6,2.2),(2.1,1.7),(2.4,2.4),(1.8,1.5),(2.6,1.9),(1.4,2.6),(2.2,2.9),(1.9,2.0)}
     \fill \p circle (0.09);
  % Cluster 2
  \foreach \p in {(7.2,2.3),(7.8,2.8),(8.1,2.1),(7.5,3.1),(6.9,2.7),(8.0,3.0),(7.3,1.8),(7.6,2.5)}
     \fill \p circle (0.09);
  % Cluster 3
  \foreach \p in {(4.7,7.2),(5.3,7.8),(4.9,8.1),(5.6,7.1),(4.4,7.9),(5.1,7.4),(5.8,7.7),(4.6,7.6)}
     \fill \p circle (0.09);
  % Testpunkte
  \draw[line width=1.1pt] (4.9,4.6) ++(-0.16,-0.16) -- ++(0.32,0.32)
        (4.9,4.6) ++(-0.16,0.16) -- ++(0.32,-0.32);
  \node[scale=0.75] at (5.75,4.6) {$P_1$};
  \draw[line width=1.1pt] (9.3,8.8) ++(-0.16,-0.16) -- ++(0.32,0.32)
        (9.3,8.8) ++(-0.16,0.16) -- ++(0.32,-0.32);
  \node[scale=0.75] at (8.6,8.8) {$P_2$};
\end{tikzpicture}
\end{center}

\textbf{a)} Nennen Sie \textbf{zwei geeignete} und \textbf{ein ungeeignetes} Verfahren zur
Novelty Detection für diese Datenlage. Begründen Sie \textbf{jedes der drei} einzeln aus der
Form der Daten -- nicht aus Laufzeit oder Implementierungsaufwand.
\antwortlinien{6}

\textbf{b)} Wie bewerten Ihre drei Verfahren jeweils den Punkt $P_1$ und den Punkt $P_2$?
Genau ein Punkt trennt die geeigneten von dem ungeeigneten Verfahren -- welcher, und warum?
\antwortlinien{4}

\textbf{c)} Welcher Metaparameter Ihres ersten geeigneten Verfahrens entscheidet darüber, ob die
drei Gruppen als \textbf{drei getrennte Inseln} oder als \textbf{eine große Region} beschrieben
werden? Was passiert bei einer schlechten Wahl in beide Richtungen?
\antwortlinien{4}

\textbf{d)} Der Datensatz wächst auf 200 Merkmale und $10^6$ Punkte. Welches Verfahren würden Sie
jetzt wählen und welche zwei Gründe sprechen gegen Ihre bisherige Wahl?
\antwortlinien{3}

\newpage
\section*{Aufgabe 8 -- Kurzfragen quer durch die Vorlesung \pkt{10}}

\textbf{a)} Beschreiben Sie in eigenen Worten, wie der \textbf{AnoGAN}-Score für ein neues Bild
berechnet wird. Warum ist dafür eine \emph{Optimierung} nötig, und welches Problem löst
\textbf{f-AnoGAN} an genau dieser Stelle?
\antwortlinien{5}

\textbf{b)} Nennen Sie die \textbf{zwei} Diversitätsquellen eines RandNet-Ensembles und
begründen Sie, warum der Ensemble-Score als \textbf{Median} und nicht als Mittelwert gebildet
wird.
\antwortlinien{4}

\textbf{c)} Vergleichen Sie GOAD und Deep SVDD in vier Zeilen:

\begin{center}
\begin{tabular}{|p{4.2cm}|p{5.2cm}|p{5.2cm}|}
\hline
 & \textbf{Deep SVDD} & \textbf{GOAD}\\ \hline
Was bilden die Normaldaten im Latent Space? & \rule{0pt}{22pt} & \\ \hline
Woher kommt das Lernsignal? & \rule{0pt}{22pt} & \\ \hline
Was ist die triviale Lösung, und was verhindert sie? & \rule{0pt}{28pt} & \\ \hline
Wie entsteht der Anomalie-Score? & \rule{0pt}{28pt} & \\ \hline
\end{tabular}
\end{center}

\vspace{6pt}
\noindent\rule{\linewidth}{0.8pt}
\begin{center}\small
\textbf{Punkteverteilung:} A1 21 $\cdot$ A2 12 $\cdot$ A3 18 $\cdot$ A4 16 $\cdot$ A5 14 $\cdot$
A6 22 $\cdot$ A7 12 $\cdot$ A8 10 \quad = \quad \textbf{125 Punkte}\\[3pt]
Lösungen liegen bewusst nicht bei -- Antworten zur Korrektur mit Begründungs-Check einreichen.
\end{center}

\end{document}
